% =========================================================
% settings.tex — Preámbulo: paquetes, formato y comandos propios
% ---------------------------------------------------------
% TODOS los \usepackage del documento van aquí; main.tex no carga
% paquetes. Junto a cada paquete se indica en qué secciones se usa,
% para saber qué se rompe si se quita o se cambia.
%
% Reglas de orden (no cambiarlas sin probar):
%   - babel antes de csquotes.
%   - xcolor[table] antes de cualquier otro paquete que cargue xcolor.
%   - float antes de caption.
%   - hyperref al final de los paquetes; las definiciones propias, después.
%   - Los tipos de columna con >{...} se definen aquí y no dentro de los
%     capítulos, porque babel (spanish) activa < y > en el cuerpo.
% =========================================================


% ---------------------------------------------------------
% 1. CODIFICACIÓN E IDIOMA                       (todo el documento)
% ---------------------------------------------------------
\usepackage[utf8]{inputenc}     % Entrada en UTF-8 (acentos, ñ)
\usepackage[T1]{fontenc}        % Codificación de fuente de salida (guiones y acentos correctos)
\usepackage[spanish]{babel}     % Idioma: separación silábica, "Capítulo", "Cuadro", etc.
\usepackage{csquotes}           % \enquote{}: 2.1 SafetiPin, 2.2 RTMDX, 2.6 Comparativa


% ---------------------------------------------------------
% 2. PÁGINA Y MÁRGENES                           (todo el documento)
% ---------------------------------------------------------
\usepackage{geometry}
\geometry{
    a4paper,
    top=2.5cm,
    bottom=3cm,
    left=3cm,
    right=3cm
}


% ---------------------------------------------------------
% 3. MATEMÁTICAS Y ALGORITMOS
% ---------------------------------------------------------
\usepackage{amsmath}            % Entornos matemáticos (align*): 5.5.3 Factibilidad económica
\usepackage{algorithm}          % Entorno flotante "algorithm": 2.3 Mobile Safe-Route (CDMX)
\usepackage{algorithmic}        % Pseudocódigo (\REQUIRE, \STATE...): 2.3 Mobile Safe-Route (CDMX)
\usepackage{xfp}                % Cálculos en compilación (\fpeval). Sin uso actual


% ---------------------------------------------------------
% 4. FIGURAS Y FLOTANTES
% ---------------------------------------------------------
\usepackage{graphicx}           % \includegraphics: portada, 3.1, 3.2, 3.3, 5.5.2
                                % \resizebox: 5.5.2, 5.5.3
\graphicspath{{./figures/}}     % Carpeta de imágenes: figures/seccionN_X/
\usepackage{float}              % Opción [H] para fijar flotantes: 1.1, 5.5.2, 5.5.3
\usepackage{caption}            % Personalización de leyendas (captions)
\usepackage{pdflscape}          % Páginas horizontales (landscape): 5.4 Análisis de riesgos


% ---------------------------------------------------------
% 5. TABLAS
% ---------------------------------------------------------
\usepackage{array}              % Columnas >{...} y \extrarowheight: 5.3, 5.5.2, columnas C/K/Q, \tablaRF
\usepackage{longtable}          % Tablas que se parten entre páginas: 5.1.2, 5.2, 5.4
\usepackage{booktabs}           % \toprule, \midrule, \bottomrule: 3.2, 5.1.2, 5.2
\usepackage[table]{xcolor}      % Colores y \cellcolor en tablas: 5.4 (semáforo de riesgos)
\usepackage{tabularx}           % Columnas X de ancho automático: 2.6, 5.5.1, \tablaRF (5.1.1)
\usepackage{multirow}           % Celdas que abarcan varias filas: 5.5.2
\usepackage{makecell}           % Encabezados de varias líneas (\makecell): 5.5.2


% ---------------------------------------------------------
% 6. TEXTO EN COLUMNAS Y TEXTO DE RELLENO
% ---------------------------------------------------------
\usepackage{multicol}           % Texto en varias columnas. Sin uso actual
\usepackage{lipsum}             % Texto de relleno. Sin uso actual; quitar en la versión final


% ---------------------------------------------------------
% 7. TÍTULOS Y ESTRUCTURA                        (todo el documento)
% ---------------------------------------------------------
\usepackage{titling}            % Personalización de \maketitle. La portada no lo usa:
                                % \thetitle viene de setup/metadata.tex
\usepackage{titlesec}           % Formato de \chapter (ver abajo)

% Formato de \chapter: título centrado, en negritas y mayúsculas
\titleformat{\chapter}[display]
  {\normalfont\Large\bfseries\centering}  % Formato general
  {\chaptertitlename\ \thechapter}        % Etiqueta: "Capítulo N"
  {20pt}                                  % Separación entre etiqueta y título
  {\Large \MakeUppercase}                 % Formato del título (mayúsculas)

% Espaciado de \chapter: \titlespacing*{comando}{izq}{antes}{después}
\titlespacing*{\chapter}{0pt}{-10pt}{20pt}


% ---------------------------------------------------------
% 8. BIBLIOGRAFÍA Y ENLACES                      (todo el documento)
% ---------------------------------------------------------
\usepackage{cite}               % Citas numéricas ordenadas y comprimidas: [1-3]
\usepackage{hyperref}           % Hipervínculos; \texorpdfstring: 6.5; \phantomsection: main.tex
                                % Debe ser el último paquete.


% =========================================================
% DEFINICIONES PROPIAS
% (van después de cargar todos los paquetes)
% =========================================================

% ---------------------------------------------------------
% A. COLORES                                     (5.4 Análisis de riesgos)
% ---------------------------------------------------------
% Semáforo de la matriz de riesgos
\definecolor{riskRed}{HTML}{C00000}
\definecolor{riskOrange}{HTML}{FFC000}
\definecolor{riskYellow}{HTML}{FFFF00}


% ---------------------------------------------------------
% B. TIPOS DE COLUMNA Y AJUSTES DE TABLAS
% ---------------------------------------------------------
% C: columna X con texto centrado                (2.6 Comparativa, 5.5.1 Factibilidad técnica)
\newcolumntype{C}{>{\centering\arraybackslash}X}

% K: columna de identificador en negritas (1.6 cm)
% Q: columna de texto que ocupa el resto del ancho
% Uso: \begin{longtable}{@{}KQ@{}}               (5.1.2 Requerimientos no funcionales, 5.2 Reglas del negocio)
\newcolumntype{K}{>{\bfseries\raggedright\arraybackslash}p{1.6cm}}
\newcolumntype{Q}{p{\dimexpr\textwidth-1.6cm-2\tabcolsep\relax}}

% Las longtable quedan alineadas al margen izquierdo y derecho, sin sangría
% (5.1.2, 5.2, 5.4)
\setlength{\LTleft}{0pt}
\setlength{\LTright}{0pt}


% ---------------------------------------------------------
% C. COMANDOS DE TEXTO
% ---------------------------------------------------------
% \keywords{lista}: línea "Palabras clave" del Resumen (documents/preliminares/resumen.tex)
\newcommand{\keywords}[1]{\textbf{Palabras clave:} #1}


% ---------------------------------------------------------
% D. FICHAS DE REQUERIMIENTOS FUNCIONALES        (5.1.1 Requerimientos funcionales)
% ---------------------------------------------------------
% Se definen en el preámbulo porque babel (spanish) activa los
% caracteres < y > dentro del documento y rompe las columnas >{...}.

% \rfAncha{texto}
% Celda que abarca las columnas 2 a 4 de la tabla.
% Ancho: 1.00 + 0.60 + 1.45 = 3.05 (en unidades de \hsize),
% más los separadores (4\tabcolsep) y las líneas (2\arrayrulewidth)
% para que alinee con el resto de la tabla.
\newcommand{\rfAncha}[1]{%
  \multicolumn{3}{>{\hsize=\dimexpr3.05\hsize+4\tabcolsep+2\arrayrulewidth\relax\raggedright\arraybackslash}X|}{#1}}

% \tablaRF{Identificador}{Nombre}{Prioridad}{Requerimiento}
%         {Entrada}{Salida}{Descripción}{Precondición}{Postcondición}
%
% Genera la ficha de un requerimiento funcional con 4 columnas:
%   col 1 (.95): etiquetas, en negritas
%   col 2 (1.00): valores
%   col 3 (.60): etiquetas, en negritas
%   col 4 (1.45): valores
\newcommand{\tablaRF}[9]{%
  \par\noindent
  \begingroup
  \setlength{\extrarowheight}{2pt}%   % Un poco de aire vertical en cada fila
  \begin{tabularx}{\linewidth}{%
      |>{\hsize=.95\hsize\raggedright\arraybackslash\bfseries}X%
      |>{\hsize=1.00\hsize\raggedright\arraybackslash}X%
      |>{\hsize=.60\hsize\raggedright\arraybackslash\bfseries}X%
      |>{\hsize=1.45\hsize\raggedright\arraybackslash}X|}
    \hline
    % Fila 1: requerimiento (celda ancha)
    Requerimiento & \rfAncha{#4} \\ \hline
    % Fila 2: identificador y nombre
    Identificador & #1 & Nombre & #2 \\ \cline{1-2}
    % Fila 3: prioridad (las dos últimas celdas quedan vacías)
    Prioridad de desarrollo & #3 & & \\ \hline
    % Fila 4: entrada y salida
    Entrada & #5 & Salida & #6 \\ \hline
    % Filas 5-7: texto largo en celda ancha
    Descripción & \rfAncha{#7} \\ \hline
    Precondición & \rfAncha{#8} \\ \hline
    Postcondición & \rfAncha{#9} \\ \hline
  \end{tabularx}%
  \endgroup
  \par\bigskip                         % Espacio antes de la siguiente tabla
}
